\section{Discussion and Final Conclusion}
\label{sec:final_discussion_conclusion}

This thesis asked which repository-derived indicators are sufficiently stable, interpretable, resistant to manipulation,
and technically feasible to support a developer-reputation proof-of-concept. No single metric dominates every criterion.
The evidence instead supports a layered result: M1 is the strongest size/context signal, M2 records change
magnitude, M3 provides the clearest activity-continuity signal, and M4 provides the best overall trade-off for the baseline
reputation-oriented PoC. Phase III then demonstrates how the selected M4 value can be made deterministic, replayable, and
auditable across the off-chain/on-chain boundary.

\subsection{Synthesis of the evidence}
\label{sec:final_evidence_synthesis}

The study combines four forms of evidence that should not be treated as interchangeable. First, Phase I provides empirical
repository measurements from ten curated projects, 80 frozen tags, and 70 tag-chronological windows. Second, C1 calculates
stability from those observations. Third, C2--C4 and the feasibility rubric encode transparent structured author assessments;
their reproducibility preserves the stated judgments but does not convert them into measured facts or independent ratings.
Fourth, Phase III supplies implementation evidence through deterministic replay, contract tests, authenticated network
transactions, and independent inspection of committed channel data.

This separation is necessary for interpreting the final decision. The 60/40 readiness weighting orders candidates within the
declared thesis framework; it is not an inferential model. Likewise, a valid Fabric transaction demonstrates that the
implemented authorization, arithmetic, endorsement, and state-transition rules operated as specified. It does not establish
that churn is semantic quality or that an off-chain claim becomes true merely because it is committed to a ledger.

\begin{table}[H]
\centering
\small
\renewcommand{\arraystretch}{1.08}
\begin{tabularx}{\textwidth}{p{1.1cm}p{4.5cm}X}
\toprule
\textbf{RQ} & \textbf{Most relevant evidence} & \textbf{Supported answer} \\
\midrule
RQ1 & Window CV, adjacent-window rank correlation, contributor overlap, short-window sensitivity, and ancestry-only sensitivity & M1 is the most stable context signal. Among behavior-oriented window signals, M4 concentration is least volatile; M3 has the strongest rank continuity among recurring developers. \\
RQ2 & C2 and C3 structured rubrics, bot sensitivity, identity policy, and retained supporting fields & M3 is easiest to explain; M4 has the strongest assessed gaming-resistance profile, but neither is a direct quality measure and the gaming result is not an attack experiment. \\
RQ3 & C1--C4 comparison, feasibility gate, deterministic replay, local contract tests, and live Fabric validation & M4 is the strongest eligible baseline PoC candidate under the declared framework and is implementable as replayable off-chain evidence plus a compact, validated ledger commitment. \\
\bottomrule
\end{tabularx}
\caption{Thesis-level answers and their evidence boundaries.}
\label{tab:final_rq_summary}
\end{table}

\subsection{Answer to RQ1: stability and comparability}
\label{sec:final_rq1}

RQ1 asks which indicators are stable across versions and comparable across projects. At window level, M1 endpoint KLOC is
the most stable signal, with primary median CV 0.014, because project size changes gradually between releases. This makes M1
valuable for context and normalization, but not as standalone developer-reputation evidence. Among behavior-oriented signals,
M4 is the more stable family: primary median CV is 0.236 for top-one churn share and 0.329 for churn Gini, compared with
0.608 for M3 commits/day and 0.890 for M2 churn/KLOC/day.

The post-hoc ancestry result materially qualifies the original analysis. Twelve of the 70 tag-chronological windows are not
linear ancestor-to-descendant intervals and contain 41.77\% of human commits and 71.55\% of human churn. Excluding them leaves
58 eligible windows. The qualitative behavioral ordering is retained: ancestry-only primary median CV is 0.247 for M4
top-one share, 0.351 for M4 Gini, 0.631 for M3, and 0.890 for M2. This supports retaining M4 for the engineering PoC, but the
reduced eight-project primary summary cannot be presented as a corrected ten-project population.

Developer-rank stability adds a different perspective. Among developers appearing in adjacent windows, M3 commits/day has
the strongest primary median Spearman correlation (0.836), while M2 churn/day and M4 churn share both have 0.608. M2 and M4
produce the same within-window ordering because dividing every developer's churn by the same window denominator changes the
unit, not the rank. Moreover, the median common-developer share is only 0.172. The supported RQ1 answer is therefore
conditional: M4 concentration is the least volatile behavior-oriented window signal, whereas M3 is the stronger continuity
signal for the small recurring-developer subset. No result justifies treating one developer score as universally stable
across changing contributor populations.

\subsection{Answer to RQ2: interpretation and manipulation}
\label{sec:final_rq2}

RQ2 asks which indicators are understandable and robust against manipulation. Under C2, M3 receives 9/10 and is the clearest
implemented reputation-adjacent signal. M1 and M2 each receive 8/10, while M4 receives 7/10 because ownership shares,
concentration, canonical identity, and the denominator must all be explained together. Interpretability alone therefore
would favor M3, not M4.

The structured C3 assessment reverses that ordering: M4 receives 9/10, M2 7/10, M3 6/10, and M1 5/10. Sustained relative
ownership was judged more difficult to manufacture than commit counts, which can be increased through splitting or
automation. This is a reasoned rubric result, not empirical proof of gaming resistance. M4 can still be distorted through
meaningless high-churn edits, dominance by one unusually large change, bot misclassification, email aliases, or spoofed
identity metadata. The implementation improves detectability by retaining numerator, denominator, bot context, window
commitments, and replayable Git evidence; it does not automatically decide whether the edits are valuable.

The supported RQ2 answer is consequently a trade-off. M3 is the easiest to communicate but needs commit-splitting and bot
safeguards. M4 is more demanding to explain but offers the strongest assessed reputation meaning and manipulation cost when
its identity and evidence rules are explicit. M1, M2, and M3 should remain supporting context rather than being discarded:
size frames the project, churn exposes the absolute work behind a share, and cadence distinguishes sustained activity from an
isolated ownership spike.

\subsection{Answer to RQ3: PoC selection and technical feasibility}
\label{sec:final_rq3}

RQ3 asks which indicator is the strongest PoC candidate under implementation constraints. M4 receives a C1--C4 comparison
average of 8.125 and feasibility 16/20 (normalized to 8.0). Its weighted readiness score is 8.075, above eligible alternatives
M3 at 7.750 and M2 at 7.100. M1's numerical readiness is high, but it is ineligible as a standalone reputation metric because
it describes project context. M5 is excluded because the current dataset contains no issue or pull-request event stream from
which responsiveness can be replayed.

The implementation results support the feasibility part of this answer. Exact integer arithmetic independently checks all
1,063 developer rows and 70 window summaries with zero failures. For the eligible \texttt{urllib3\_W03} fixture, independent
Git extraction reproduces numerator 325, denominator 1,058, and 307,183 ppm. A standalone bundle supports fresh-repository
replay; thirteen off-chain tests and an independent Git traversal confirm the extraction and commitment chain. The unchanged
JavaScript contract then passes 34 local tests covering authorization, schema constraints, \texttt{BigInt} arithmetic,
immutability, indexes, history, and events.

The live network closes the remaining implementation gap. The package was installed on both peers and approved by both
organizations, and the
committed validation parameter was independently decoded as requiring both organizational peers. All 17 authenticated Gateway
checks and all 47 independent network checks pass. The accepted Org1 transaction is committed with validation code 0 in block
6, is read identically through both peers, appears in both indexes and history, and emits the expected event. Malformed JSON,
a one-ppm mutation, an Org2 submission, a duplicate, and a temporary one-byte artifact modification are rejected through their
respective on-chain or replay boundaries.

The supported RQ3 answer is that M4 is feasible in a hybrid architecture, not that Git analysis should occur in chaincode.
Git traversal, ancestry, bot classification, identity normalization, and aggregation remain off-chain; Fabric authenticates
the configured submitter, recomputes the compact arithmetic, requires the declared endorsements, prevents overwrite, and binds
the result to replayable evidence through hashes.

\subsection{Implications for a DESTINO-aligned reputation system}
\label{sec:final_implications}

The main architectural implication is that a reputation system should preserve evidence layers instead of reducing activity
to an opaque score. A compact ledger record can identify the project, window, pseudonymous subject, policy version, exact
integer value, numerator, denominator, and evidence commitments. An auditor can then retrieve the matching artifact and
reapply the same policy. This design makes disagreement inspectable: a party can challenge the artifact, extraction,
identity rule, or interpretation without changing the already committed record.

Policy versioning is equally important. Changing bot rules, identity linkage, release-window eligibility, or rounding creates
a different measurement contract and should produce a new policy version rather than silently rewriting prior evidence. The
ledger's role is therefore provenance, authorization, ordering, and immutability under configured governance. It is not an
oracle for developer identity, semantic code quality, or the long-term availability of external artifacts.

\subsection{Limitations and future work}
\label{sec:final_future_work}

The limitations in Section~\ref{sec:threats_validity} bound the conclusion. The ten repositories are curated rather than a
statistical sample of all software ecosystems. The ancestry-only subset removes all JUnit windows and leaves only eight
projects able to contribute to its primary stability summary. Email canonicalization is pseudonymous and does not resolve
aliases or prove real-world identity. C2--C4 are structured author assessments without independent raters, and the 60/40
weighting is a design choice. Finally, the Fabric evaluation uses one fixture, generated test identities, one peer per
organization, one orderer, and no performance or fault experiment.

The most valuable extensions would be to replay multiple eligible windows and developers, add issue and pull-request evidence
for M5, evaluate explicit identity-linking governance, recruit independent raters for C2--C4, and conduct controlled gaming
experiments. A production-oriented systems evaluation would additionally require multiple hosts, richer certificate
governance, more peers and orderers, privacy controls, load measurements, failures, backup and recovery, and an externally
durable artifact store. These extensions would broaden external validity; they are not required to change the completed
baseline claim.

\subsection{Final conclusion}
\label{sec:final_conclusion}

This thesis contributes a reproducible path from repository history to auditable ledger evidence. It constructs and validates
a multi-project dataset, compares candidate indicators under explicit criteria, identifies and discloses a material
release-window ancestry limitation, selects M4 without changing the frozen evidence, and implements the selected signal through deterministic
off-chain verification and Hyperledger Fabric chaincode.

Relative to the broad metric and assessment literature and the full DESTINO proposal, this contribution is intentionally
focused. It complements the earlier general perspective by operationalizing one selected repository signal as an auditable
evidence record; it does not claim to provide a universal developer-quality measure, a complete reputation model, or the
complete DESTINO ecosystem.

The final conclusion is deliberately bounded. Under the declared evaluation framework and ancestry-only eligibility gate,
M4 developer churn ownership share is a suitable signal for the baseline DESTINO-aligned PoC. The implemented mechanism
reproduces the selected fixture from Git, validates exact arithmetic and evidence commitments, and commits the resulting record
through authenticated two-organization endorsement with independently confirmed channel state. This establishes a coherent,
replayable, and testable PoC. It does not establish that M4 universally measures developer quality, that gaming resistance is
empirically proven, or that the educational Fabric topology is production-ready. Within those limits, the research questions
are answered and the baseline thesis objective is achieved.
